\documentclass{article}
\usepackage[utf8]{inputenc}
\usepackage[T1]{fontenc}
\usepackage{geometry}
\usepackage{graphicx}
\usepackage{float}
\usepackage{hyperref}
\usepackage[english,indonesian]{babel}

\title{Optimasi Pemrosesan Template LaTeX Menggunakan Pendekatan Berbasis AI}
\author{Academic Support Agent}
\date{\today}

\begin{document}
\maketitle

\begin{abstract}

\end{abstract}

\chapter{PENDAHULUAN}


\chapter{TINJAUAN PUSTAKA}


% Original content (commented out):
% \documentclass[11pt]{article}

\usepackage[utf8x]{inputenc}
\usepackage{amsmath}
\usepackage{graphicx,geometry}

\newgeometry{margin=2cm}

\title{Your Paper}
\author{You}
\newlength{\thelinewidth}

\thelinewidth=\textwidth

\newlength{\exmpwidinf}
\newlength{\exmpwidouf}
\newlength{\exmpwidewid}

\exmpwidinf=0.43\thelinewidth
\exmpwidouf=0.43\thelinewidth
\exmpwidewid=0.9\thelinewidth

\newenvironment{example2}{
    \ttfamily\obeylines\obeyspaces\frenchspacing
}

\newenvironment{example}[1][]{
    \ttfamily\obeylines\obeyspaces\frenchspacing
    \newcommand{\exmp}[2]{
        \begin{minipage}[t]{\exmpwidinf}\rightskip=0pt plus 1fill\relax##1\medskip\end{minipage}&
        \begin{minipage}[t]{\exmpwidouf}\rightskip=0pt plus 1fill\relax##2\medskip\end{minipage}\\
        \hline
    }

    \begin{tabular}{|l|l|}
        \hline
        \multicolumn{1}{|c|}{\bf\texttt{superbugrahan.gir}}&
        \multicolumn{1}{|c|}{\bf\texttt{superbugrahan.cik}}\\
        \hline
}{
    \end{tabular}
}


\begin{document}

\section*{\centering Süper Buğrahan}

	
Elinizde NxN lik bir harita ve süper bir Buğrahan var.
\begin{example2}
\noindent.......
\noindent..2.2..
\noindent.21112.
\noindent..1B1..
\noindent.21112.
\noindent..2.2..
\noindent.......
\end{example2}
\noindent Haritadaki B Buğrahanı temsil ediyor, Buğrahanın hareket yetenekleri ise şu şekilde 8 yöndeki komşularından birine geçebiliyor(Haritada 1 ile gösterilen yerler) veya adeta bir at gibi L çizerek atlayabiliyor(Haritada 2 ile gösterilen yerler) ama haritanın dışına çıkamıyor.
\\\\
Sizin yapmanız gereken ise Buğrahanın NxN lik haritada Başlangıç noktasından Bitiş noktasına tam olarak M hamlede gitmesinin kaç farklı şekilde mümkün olacağını hesaplayıp Buğrahana yardım etmek. Koordinatlar 0'dan N-1'e gidecek şekilde tutulmaktadır.

\subsection*{Girdi}

İlk satırda Haritanın boyutunu belirten $3 \leq N \leq 100$.
\\\\
\noindent İkinci satırda Başlangıç noktasının koordinatları(ilki satır ikincisi sütun numarası olacak şekilde)
\\\\
\noindent Üçüncü satırda Bitiş noktasının koordinatları(ilki satır ikincisi sütun numarası olacak şekilde)
\\\\
\noindent Başlangıç ve Bitiş noktalarının haritanın içinde olacağı garanti edilmektedir.
\\\\
\noindent Ve son satırda istenilen hareket sayısını belirten $1 \leq M \leq 50$.

\subsection*{Çıktı}

\noindent Buğrahanın Başlangıç noktasından Bitiş noktasına tam olarak M hamlede kaç farklı şekilde gidebileceğini belirten sayı. (Bu sayının Long Long inte sığacağı garanti edilmektedir)

\subsection*{Örnek}

\begin{example}%
\exmp{
3
0 0
1 0
1
}{
1
}%
\exmp{
3
0 0
2 2
1
}{
0
}%
\exmp{
3
0 0
0 0
2
}{
5
}%
\end{example}
   

\end{document}


\end{document}